\section{总结与延伸}

本讲从 base model 的 control gap 出发，构建了完整的 mid/post-training 视角。SFT 用 demonstrations 教会模型 interaction protocol；mid-training 用更大规模高质量行为数据重塑能力；preference data 把 generation 转为 comparison；PPO 通过 on-policy RL 优化 reward；DPO 则把受约束 RLHF 转成离线 pairwise loss。

六个关键结论是：

\begin{enumerate}
\item SFT 最擅长提取和组织已有能力，不是可靠的长尾知识写入器。
\item Style、length 和 formatting 会显著影响 preference eval，必须作为 confounders 审计。
\item Preference labels 取决于 annotator/judge 分布，不是客观 reward。
\item PPO 强大但系统复杂；DPO 简洁但受 offline support 限制。
\item 所有方法都可能 overoptimize proxy reward，导致 mode collapse 和 calibration 退化。
\item Data recipe、sampling、base model 和 evaluation 常与算法本身同等重要。
\end{enumerate}

下一讲将 reward 从“人类觉得哪个好”换成“答案能否被自动验证”，讨论 PPO、GRPO、DeepSeek R1、Kimi K1.5、Qwen3 以及 RLVR 的系统问题。

\subsection{拓展阅读}

下面的阅读路线按“公开 pipeline→偏好算法→数据与失败模式”排序。阅读时不要只抄最终 loss；应同时记录数据由谁生成、candidate 来自哪个 policy、judge 怎样定义赢家、优化是否在线，以及论文用什么独立指标证明 reward 上升没有变成 Goodhart failure。

\begin{longtable}{L{0.23\textwidth}L{0.31\textwidth}L{0.36\textwidth}}
\toprule
材料 & 为什么值得读 & 带着什么问题读 \\
\midrule
Stiennon et al., \emph{Learning to Summarize from Human Feedback} & 早期公开材料中对 preference collection、reward model 与 PPO 流程披露较完整 & 标注指南怎样把“好摘要”拆成可比较维度？ \\
Ouyang et al., \emph{Training Language Models to Follow Instructions with Human Feedback} & InstructGPT 的 SFT→RM→PPO 三阶段主线，也是本讲多数系统图的来源 & 1.3B/6B/175B 模型的 data、KL 与 human eval 如何共同作用？ \\
Bai et al., \emph{Constitutional AI} & 展示 critique/revision、自训练与 AI feedback 如何替代部分人工标签 & Constitution 把哪些价值判断显式化，又留下哪些 judge bias？ \\
Rafailov et al., \emph{Direct Preference Optimization} & 从 KL-regularized RLHF 的最优 policy 推导离线 pairwise objective & Nonparametric assumption、reference policy 与 $\beta$ 分别承担什么角色？ \\
Tulu 3 技术报告 & 提供较现代的开源 post-training mixture、偏好数据与 DPO 变体实践 & 哪些收益来自算法，哪些来自 candidate generation、过滤与 verifier？ \\
Santurkar et al. 与 Hosking et al. 的 annotator 研究 & 直接展示 demographics、expertise 与 style preference 如何改变聚合标签 & “更多标注者”何时只降低方差，却无法修复 population bias？ \\
Reward overoptimization 与 length-bias 研究 & 揭示 proxy reward 先升后降、verbosity shortcut 与 noiseless/noisy feedback 的区别 & 转折点如何被 fresh human eval、KL、length 与 entropy 指标提前发现？ \\
\bottomrule
\caption{Lecture 15 的注释式拓展阅读路线。}
\end{longtable}

\begin{knowledgebox}{建议的复现实验}
从一个小型 instruction model 出发，构造长度平衡的 pairwise dataset，分别训练 winner-only SFT 与 DPO；在 held-out pairs 上同时报告普通 win rate、length-controlled win rate、事实正确率、平均长度和输出 entropy。这个实验能把本讲的数据偏差、离线支持集和过度优化问题放进同一张结果表。
\end{knowledgebox}

\end{document}
